% !TEX root = ../notes.tex

\documentclass[../notes.tex]{subfiles}

\begin{document}

\section{May 6}
Here we go.

\subsection{Deformations}
Here is an example of a central extension.
\begin{example}
	Let $\mf g$ be the abelian Lie algebra $k^2$ with basis $\{x,y\}$. Then $\mathrm H^\bullet(\mf g;k)=\land^\bullet k^{2*}$. Thus, $\mathrm H^2(\mf g;k)$ is one-dimensional, so there is a nontrivial extension given by $\omega$ with $\omega(x,y)=1$. The corresponding Lie algebra $\widetilde{\mf g}$ is then a central abelian extension
	\[0\to kz\to\widehat{\mf g}\to kx\oplus ky\to0\]
	satisfying $[x,y]=z$ and $[x,z]=[y,z]=0$. This $\widetilde{\mf g}$ can be seen to be the Heisenberg Lie algebra.
\end{example}
One can try to do more with formal deformations, such as all the way to $k[[t]]$.
\begin{proposition} \label{prop:deform-lie-algebra}
	Let $\mf g$ be a Lie algebra over a field $k$. If
	\[\mathrm H^2(\mf g;\mf g)=0,\]
	then $\mf g$ admits a unique formal deformation to $k[[t]]$.
\end{proposition}
\begin{proof}
	We are asking for a Lie algebra $\widetilde{\mf g}$ over $k[[t]]$, which becomes $\mf g$ upon sending $t\to0$. This means that we are asking for a commutator
	\[[x,y]=c_0(x,y)+tc_1(x,y)+t^2c_2(x,y)+\cdots,\]
	where $c_0(x,y)$ is the commutator of $\mf g$, and each $c_i$ is a map $\land^2\mf g\to\mf g$. Now, one can check that $[-,-]$ satisfying the Jacobi identity is equivalent to
	\[dc_n=\sum_{1\le i\le n-1}[c_i,c_{n-i}],\]
	which is called the Schouten bracket.
	\begin{example}
		One has $dc_2=[c_1,c_1]$, meaning that
		\[dc_2(x,y,z)=c_1(c_1(x,y),z)+c_1(c_1(y,z),x)+c_1(c_1(z,x),y).\]
	\end{example}
	Now, one can check that
	\[d\Bigg(\sum_{1\le i\le n-1}[c_i,c_{n-i}]\Bigg)=0,\]
	so we can find a deformation to $c_n$ if and only if the corresponding cocycle $\gamma\in\mathrm H^3(\mf g;\mf g)$ vanishes. In fact, one can see that the choice of $c_n$ is unique up to an element of $\mathrm H^2(\mf g;\mf g)$. The vanishing of these groups proves the theorem.

	Now, given that $\mathrm H^2(\mf g;\mf g)=0$, we see that we may change coordinates to force $c_1=0$. Then the class in $\mathrm H^3(\mf g;\mf g)$ vanishes, and we can lift to $c_2$. This $c_2$ is unique, but we note that we could also lift $c_2=0$, so a change of coordinates must be able to force $c_2=0$. This argument continues inductively.
\end{proof}
\begin{example}
	If $\mf g$ is semisimple, then $\mathrm H^2(\mf g;\mf g)$ vanishes by \Cref{cor:semisimple-has-vanishing-cohomology}, so we conclude that deformations are trivial by \Cref{prop:deform-lie-algebra}.
\end{example}
\begin{remark}
	In fact, if merely $\mathrm H^2(\mf g;\mf g)=0$, then we see that we are forced to take $c_1=0$, so the class in $\mathrm H^3(\mf g;\mf g)$ automatically vanishes.
\end{remark}
\begin{remark}
	If $\mathrm H^3(\mf g;\mf g)=0$, then the argument shows that we can always extend the deformation, but we always have many choices given by $\mathrm H^2(\mf g;\mf g)$.
\end{remark}
\begin{remark}
	If $\mathrm H^2(\mf g;\mf g)$ and $\mathrm H^3(\mf g;\mf g)$ are both nonzero, then understanding deformations is rather tricky because the choice of obstruction at level $c_n$ depends on the choices for $\{c_1,\ldots,c_{n-1}\}$. It is still possible to build deformations in some special cases.
\end{remark}
\begin{example}
	Suppose that $\mf g$ is abelian. Then the equation for $n=2$ is $[c_1,c_1]=0$. Notably, $c_1\in\mathrm H^2(\mf g;\mf g)$ has equivalent data to a map $\land^2\mf g\to\mf g$, and it always provides a deformation of $\mf g$ to $k[t]/\left(t^2\right)$, but they only lift to $k[t]/\left(t^3\right)$ if $[c_1,c_1]=0$. However, once $[c_1,c_1]=0$ so that we lift to $k[t]/\left(t^3\right)$, then we can lift all the way to $k[[t]]$.
\end{example}
\begin{remark}
	We may say that obstructions to deforming the Lie algebra $\mf g$ live in $\mathrm H^3(\mf g;\mf g)$ with freedom $\mathrm H^2(\mf g;\mf g)$. This is a rather general notion. For example, given a $\mf g$-module $V$, we may ask if we can deform to $k[[t]]$. One can show using similar methods that deformations to $k[t]/\left(t^2\right)$ are classified by $\op{Ext}^1_{\mf g}(V,V)=\mathrm H^1(\mf g;\op{End}_kV)$, and obstructions to lifting further are in $\mathrm H^2(\mf g;\op{End}_kV)$.
\end{remark}

\subsection{The Levi Decomposition}
As an application of our cohomological discussion, let's prove the Levi decomposition.
\begin{theorem}[Levi decomposition] \label{thm:levi}
	Fix a finite-dimensional complex Lie algebra $\mf g$, and let $\op{rad}\mf g$ be the largest solvable ideal. Then the short exact sequence
	\[0\to\op{rad}\mf g\to\mf g\to\mf g_{\mathrm{ss}}\to0\]
	splits.
\end{theorem}
\begin{proof}
	We begin by choosing some splitting $\mf g=\mf g_{\mathrm{ss}}\oplus\rad\mf g$ of vector spaces, which we use to write down elements of $\mf g$. Then the commutator looks like
	\[[(a,x),(b,y)]=([x,b]-[y,a]+[a,b]+\omega(x,y),[x,y]),\]
	for some $\omega\colon\land^2\mf g_{\mathrm{ss}}\to\mf g_{\mathrm{ss}}$. Observe that if $\omega$ vanished, then this commutator is just the semidirect product commutator $\mf g_{\mathrm{ss}}\ltimes\rad\mf g$, so we would be done. Thus, our goal is to modify the splitting in order to remove $\omega$.

	To this end, we would like to make $\rad\mf g$ abelian. This is not literally possible, but it is almost possible. Because $\mf g$ is solvable, we note that $\op{rad}\mf g$ is solvable and hence admits a filtration
	\[D^0\supseteq D^1\supseteq\cdots\subseteq D^{n+1}=0,\]where $D^n$ is an abelian ideal, and $D^{i+1}=\left[D^i,D^i\right]$. We now proceed by induction on $d=\dim\rad\mf g$, where the base case of $d=0$ has no content.

	For the induction, we may set $\overline{\mf g}\coloneqq\mf g/D^n$, and the Levi decomposition holds for $\overline{\mf g}$ by the induction. Thus, we may choose a splitting for $\overline{\mf g}$, and choosing the splitting of $\mf g$ to lift this one, we find that $\omega$ outputs to $D^n$. The moral is that $\omega$ is a $2$-cocycle in $Z^2(\mf g_{\mathrm{ss}};D^n)$. However, \Cref{cor:semisimple-has-vanishing-cohomology} tells us that $\mathrm H^2(\mf g_{\mathrm{ss}};D^n)=0$, so $\omega$ is a $2$-coboundary, so we may find some $f\colon\mf g_{\mathrm{ss}}\to D^n$ with
	\[\omega(x,y)=f([x,y])-xf(y)+yf(x).\]
	Now, modifying the choice of splitting $\mf g_{\mathrm{ss}}\to\mf g$ by adding in $f$ causes $\omega$ to disappear, and we are done by the first paragraph.
\end{proof}

\subsection{The Third Fundamental Theorem}
Let's start by recalling the first two fundamental theorems.
\begin{theorem} \label{thm:lie-1}
	Fix a Lie group $G$ with Lie algebra $\mf g$. For each Lie subalgebra $\mf h\subseteq\mf g$, there is a unique closed connected Lie subgroup $H\subseteq G$ with $\op{Lie}H=\mf h$.
\end{theorem}
\begin{theorem} \label{thm:lie-2}
	Fix Lie groups $H$ and $G$ with Lie algebras $\mf h$ and $\mf g$. If $H$ is connected, then the differential map
	\[\op{Hom}(H,G)\to\op{Hom}(\mf h,\mf g)\]
	is injective, and if $H$ is simply connected, then it is an isomorphism.
\end{theorem}
The moral of \Cref{thm:lie-2} is that the functor $\op{Lie}$ defines a fully faithful embedding from the category of simply connected Lie groups to the category of finite-dimensional Lie algebras. In fact, this functor is an equivalence, which is the assertion of the third theorem.
\begin{theorem} \label{thm:lie-3}
	Fix a finite-dimensional real Lie algebra $\mf g$. Then there is a Lie group $G$ with Lie algebra $\mf g$.
\end{theorem}
\begin{remark}
	There is an analogous statement over $\CC$, and the proof is the same.
\end{remark}
One possible way to prove this is to prove Ado's theorem, which shows that a finite-dimensional Lie algebra has a faithful representation. Then we can use \Cref{thm:lie-1} to conclude. We will instead prove \Cref{thm:lie-3} directly, using the Levi decomposition.
\begin{lemma} \label{lem:lie-3-solvable}
	\Cref{thm:lie-3} is true if $\mf g$ is solvable.
\end{lemma}
\begin{proof}
	We proceed by induction on $\dim\mf g$, where the case of $\mf g=0$ has no content. Because $\mf g$ is solvable, its filtration by ideals reveals that $\mf g$ admits a nontrivial character $\chi\colon\mf g\to\RR$. Then $\mf g_0\coloneqq\ker\chi$ has dimension smaller than $\mf g$, so the induction allows us to find a simply connected Lie group $G_0$ with $\op{Lie}G_0=\mf g_0$.

	Now, choose $d\in\mf g$ with $\chi(d)=1$, and one can see that $\mf g=\RR d\ltimes\mf g_0$ (the splitting is automatic because $\RR d$ is one-dimensional), where $d$ acts by a derivation on $\mf g_0$. Thus, $\exp(t\op{ad}d)$ gives a one-parameter subgroup of $\op{Aut}\mf g_0$, so \Cref{thm:lie-2} lifts this to a one-parameter subgroup $\{e^{td}\}$ of automorphisms of $G_0$. Accordingly, we define the group $G\coloneqq G_0\rtimes\RR$, whose operation is given by
	\[(x,t)\cdot(y,s)\coloneqq\left(x\cdot e^{td}(y),t+s\right).\]
	From this construction, one can directly check that $\op{Lie}G=\mf g$.
\end{proof}
\begin{remark} \label{rem:nilpotent-lie-3}
	Further assume that $\mf g$ is nilpotent. Then one can use this argument to show $\exp\colon\mf g\to G$ is a diffeomorphism and that the multiplication map $G\times G\to G$ is a polynomial $\mu_0$ in the coordinates of $\mf g$. By the induction, we may assume this is true for $\mf g_0$, so we may identify $G_0$ with $\mf g_0$ by $\exp$, and $G$ has a multiplication law
	\[(x,t)\cdot(y,s)=\left(\mu_0(x,e^{td}y),s+t\right).\]
	Because $\mf g$ is nilpotent, some power of the derivation $d$ vanishes, so $e^{td}$ expands out as some polynomial. Thus, the multiplication map $G\times G\to G$ is polynomial. As for $\exp$ being a diffeomorphism, we can compute $\exp$ as
	\[\exp(x,t)=\left(\exp\left(\frac{e^{td}-1}{td}(x)\right),t\right).\]
	(Namely, one should check that this is a group homomorphism with the correct differential.) By inverting $\frac{e^{td}-1}{td}$ (which is possible because $d$ is nilpotent), we can solve for $x$ in terms of the right-hand side, which provides the inverse of $\exp$.
\end{remark}
\begin{remark}
	One can use the idea of \Cref{rem:nilpotent-lie-3} to prove \Cref{thm:lie-3} for nilpotent Lie algebras directly by using the Campbell--Hausdorff formula in order to define a Lie group structure directly on the space $\mf g$.
\end{remark}
We are now ready to prove the general case.
\begin{proof}[Proof of \Cref{thm:lie-3}]
	By \Cref{thm:levi}, we have a splitting $\mf g=\mf g_{\mathrm{ss}}\ltimes\mf a$, where $\mf a\coloneqq\rad\mf g$. Now, we can let $G_{\mathrm{ss}}$ be the simply connected cover of $G_{\mathrm{ad}}$, which has Lie algebra $\mf g_{\mathrm{ss}}$. Similarly, \Cref{lem:lie-3-solvable} allows us to lift $\mf a$ to a Lie group $A$.

	Now, the natural map $\mf g_{\mathrm{ss}}\to\op{Der}\mf a$ lifts by \Cref{thm:lie-2} to a map $G_{\mathrm{ss}}\to\op{Aut}(A)$. This map allows us to define a semidirect product $G\coloneqq G_{\mathrm{ss}}\ltimes A$, and as in \Cref{lem:lie-3-solvable}, one finds that $\op{Lie}G=\mf g$. This completes the proof.
\end{proof}
\begin{remark}
	By the construction of \Cref{lem:lie-3-solvable}, we see that $A$ is diffeomorphic to a Euclidean space. Now, by the equivalence of categories, we see that there is a unique simply connected Lie group $G$ with Lie algebra $\mf g=\mf g_{\mathrm{ss}}\ltimes\mf a$, and by the construction of the proof of \Cref{thm:lie-3}, we see that $G$ is homotopic to $G_{\mathrm{ss}}$, which is homotopic to $G_{\mathrm{ss}}^c$ by \Cref{thm:polar-decomp}. For example, the cohomology of (simply connected) Lie groups reduces to the cohomology of compact groups!
\end{remark}

\subsection{Flag Manifolds}
The last topic of the course will be flag manifolds. Let's start by defining Borel subgroups.
\begin{notation}
	Throughout, $G$ is a connected complex reductive group with Lie algebra $\mf g$, and we let $\mf h\subseteq\mf g$ be a Cartan subalgebra with corresponding maximal torus $H\subseteq G$. We let $\Phi_+$ be a choice of positive roots, which induces a decomposition $\mf g=\mf n_-\oplus\mf h\oplus\mf n_+$; let $N_\pm\subseteq G$ be the Lie group corresponding to $\mf n_\pm$.
\end{notation}
\begin{definition}[Borel subgroup]
	Fix everything as above. Then a \textit{Borel subgroup} is one conjugate to $B_+\coloneqq HN_+$.
\end{definition}
\begin{remark}
	It is more typical to define $B$ as a maximal solvable subgroup, but this is equivalent.
\end{remark}
\begin{remark}
	Because all Cartan subalgebras $\mf h$ are conjugate, the definition of a Borel subgroup does not depend on the choice of Cartan or positive roots.
\end{remark}
Here is a sanity check, which explains why $B_+$ has a chance of being a maximal solvable subgroup.
\begin{lemma} \label{lem:borel-is-own-normalizer}
	Fix everything as above with a complex reductive group $G$. Then $N_G(B_+)=B_+$.
\end{lemma}
\begin{proof}
	Suppose that $g$ normalizes $B_+$. Note that $H'\coloneqq gHg^{-1}$ is another maximal torus, and it is a torus in $B_+$. In fact, there is $b\in B_+$ with $bH'b^{-1}=H$. Indeed, $H'$ and $H$ are already the same modulo $N_+$, so $\op{Lie}H'$ and $\op{Lie}H$ agree modulo $\mf n_+$; up to conjugating by a simple root space in $B_+$, we can force them to agree up to $[\mf n_+,\mf n_+]$, and then we can iterate the process. Thus, we may assume that $H=H'$, so $g$ is in the normalizer of $H$ while also fixing positive roots (because it normalizes $B$), so $g$ has trivial image in the Weyl group, so $g\in H$.
\end{proof}
We are now ready to define the flag manifold.
\begin{definition}[flag manifold]
	Fix everything as above with a complex reductive group $G$. Then we define $\mc F$ to be the set of all Borel subgroups.
\end{definition}
\begin{remark}
	By definition, $\mc F$ admits a natural transitive action of $G$ by conjugation. The fiber over $B_+$ is $N(B_+)$, which is $B_+$ by \Cref{lem:borel-is-own-normalizer}, so $\mc F=G/B$. Thus, $\mc F$ attains the structure of a homogeneous space and thus a complex manifold.
\end{remark}
\begin{remark}
	Adding in central torus to $G$ will not change $\mc F$ because the central torus goes straight into $B$. Thus, $\mc F$ depends only on $\mf g_{\mathrm{ss}}$.
\end{remark}
We would now like to pass to a compact group.
\begin{lemma} \label{lem:compact-on-flag}
	Fix a complex reductive group $G$, and let $G^c\subseteq G$ be the compact form. Then the action of $G^c$ on the flag manifold $\mc F$ is transitive.
\end{lemma}
\begin{proof}
	Fix a Borel subgroup $B_+$. Then the image of the map $G^c\to\mc F$ induced by $g\mapsto gB_+g^{-1}$ must be compact and hence closed. Because $\mc F$ is connected, it is enough to check that this image is also open.
	
	To this end, we claim that $B_+\in\mc F$ is an interior point of $G^c\cdot B_+$: indeed, one can check that $\mf g^c\oplus\mf b_+=\mf g$ by the direct definition of $\mf g^c$ and $\mf b_+$, we find that the tangent space map $\mf g^c\to T_{B_+}\mc F$ is surjective, so $G^c\to\mc F$ is a local diffeomorphism at $B_+$. By homogeneity of this argument, we see that the image $G^c\to\mc F$ is also open, so we are done.
\end{proof}
\begin{proposition} \label{prop:flag-as-compact}
	Fix a complex reductive group $G$ with maximal torus $H$, and let $G^c\subseteq G$ be the compact form. Then the flag manifold $\mc F$ is isomorphic to $G^c/H^c$.
\end{proposition}
\begin{proof}
	In light of \Cref{lem:compact-on-flag}, it only remains to compute the fiber of the natural map $G^c\to\mc F$. This is $G^c\cap B_+$, which we claim is $H^c$. Indeed, certainly $H^c\subseteq G^c\cap B_+$. For the converse, note that $G^c$ is $\tau$-stable and so $G^c\cap B_+$ is $w_0$-stable, so $G^c\cap B_+\subseteq w_0(B_+)\cap B_+$, which is $H$. But then $G^c\cap H=H^c$.
\end{proof}
\begin{example}
	For $G=\op{GL}_n(\CC)$, we find that $\mc F=G/B_+$ is the flag manifold $\mc F_n(\CC)$ from earlier.
\end{example}
As an application, we can prove the Iwasawa decomposition.
\begin{theorem}[Iwasawa decomposition]
	Fix a complex connected reductive group $G$. Let $K$ be the compact form, $A\coloneqq\exp(i\mf h^c)$, and $N\coloneqq N_+$. Then the multiplication map
	\[K\times A\times N\to G\]
	is a diffeomorphism.
\end{theorem}
\begin{proof}
	Consider the action of $K\times HN$ on $G$, where $K$ acts on the left, $HN$ acts on the right. By \Cref{lem:compact-on-flag}, we see that $K$ acts transitively on $G/HN$, so the action of $K\times HN$ on $G$ is transitive. Now, $K\cap HN=H^c$ by \Cref{prop:flag-as-compact}, so $K\times AN$ acts transitively on $G$ with trivial stabilizers. Thus, the induced fiber bundle projection is a diffeomorphism, as required.
\end{proof}

\end{document}